% =====================================================================
% Appendix: Storage I/O Without the Blocking Pool.
%
% Records an investigation this version leaves open: whether the
% backend's file I/O could run on the Pingora workers instead of Tokio's
% blocking pool, what that costs when the kernel holds writes back, and
% what io_uring and the choice of file system could change. Nothing here
% is normative; it is the starting point for a later version.
% =====================================================================

This appendix records an investigation that this version leaves open:
whether the backend could do its file I/O on the worker threads that poll
it, instead of handing each operation to Tokio's blocking pool as
\reqref{ADR-STORAGE.EXECUTOR-01} decides. It collects what was measured,
what the kernel allows, and the directions a later version could take.
Nothing here is normative: no requirement, decision or test case of this
version depends on it. A later version that takes the work up would
complement \reqref{ADR-STORAGE.EXECUTOR-01} and
\reqref{FDS-STORAGE.EXECUTOR-01}, as \reqref{ADR-STORAGE.PAGECACHE-01} does
for reads, and each mechanism it adds would have to meet
\reqref{HLR-SCOPE.V1-01}: improve, by at least 20\,\%, a figure that a
performance test case measures.

All the measurements were made in October 2026 on one host: an Intel Core
i7-4770HQ processor (four cores, eight hardware threads), 16~GB of memory,
and an Apple SM0256G SSD holding an ext4 file system mounted with
\texttt{relatime}, under Linux 6.12. They guide the work; they do not set
its targets. What the kernel does is cited from the Linux 6.12 source.

\section{What the Blocking Pool Still Costs}
\label{app:executor-cost}

A handoff to the blocking pool wakes a pool thread, which runs the
operation and then wakes the worker. Measured on 23~October~2026, one
handoff took 8 to 11~microseconds, from the call to its return, on an idle
runtime, and about 20~microseconds of processor time under load, shared
between the worker and the pool thread.

On Linux, the reads that the page cache answers no longer pay it
(\reqref{ADR-STORAGE.PAGECACHE-01}); stores still do. A store makes at
least three handoffs: one to create its temporary and write the record's
fixed part, one per body chunk, and one to publish it
(\reqref{FDS-STORAGE.EXECUTOR-01}). A synchronised store is dominated by its
\texttt{fdatasync}, which held a pool thread for 16 to 25~ms on this host.
Without it (\reqref{ADR-PERSISTENCE.UNSYNCED-01}), the handoffs are a large
share of what a store costs.

Two earlier attempts to cut what the pool costs failed:
\begin{itemize}
  \item \emph{Capping the pool} (\texttt{ServerConf::max\_blocking\_threads}),
    on 22~October~2026, as \reqref{ADR-STORAGE.PAGECACHE-01} reports among
    its alternatives: a cap small enough to speed hits up starved the
    synchronised stores, each of whose syncs holds a pool thread.
  \item \emph{Threads of the backend's own}, beside Tokio's pool, which the
    backend fed directly, on 23~October~2026. Without spinning, stores ran
    as fast as on Tokio's pool. Spinning for 20 to 100~microseconds on an
    empty queue, so that a thread would take the next operation before
    falling asleep, made them 3 to 9\,\% slower, since the server's
    processors were already busy. A handoff costs much the same whichever
    pool it goes to.
\end{itemize}

\section{The Experiment: File I/O on the Workers}
\label{app:executor-experiment}

The idea tested on 30~October~2026: give the server more worker threads,
and run each storage operation on the worker that polls it, with no handoff
at all. A worker blocked in a system call then stalls only the tasks it
holds, and with work stealing on, the Pingora default, the other workers
take over the tasks queued on it.

A prototype changed only the function through which every operation
reaches the pool, and the close of a removed file
(\reqref{FDS-STORAGE.EXECUTOR-01}). An environment variable, read once,
chose among three executors:
\begin{itemize}
  \item \emph{pool}: \reqref{ADR-STORAGE.EXECUTOR-01} as specified, the
    reference;
  \item \emph{inline}: each operation runs on the worker that polls it;
  \item \emph{block in place}: each operation runs on the worker too,
    inside Tokio's \texttt{block\_in\_place}, which first hands the
    worker's other tasks to another thread.
\end{itemize}
Page-cache reads (\reqref{ADR-STORAGE.PAGECACHE-01}) were on unless stated
otherwise.

\textbf{Method.} The setup of \reqref{TC-PERFORMANCE.HITRATE-01}, run by a
script of its own: the fixture proxy on four of the eight hardware threads,
the load generator and the origin (nginx) on the other four; 64 client
connections; 10~seconds per scenario; three rounds, in a different order
each round, reporting the medians. Hit scenarios first stored 100 objects
and then requested them at random, so every hit was on an object the page
cache held. Store scenarios requested a new URL every time, so every request
was a miss that the proxy fetched and stored. Each configuration started
with an empty cache. Unsynchronised stores were measured with
\texttt{Durability::NoStoreSync} (\reqref{ADR-PERSISTENCE.UNSYNCED-01}),
synchronised ones with the default.

\textbf{Hits} ran within 5\,\% of the reference with every executor and
every number of workers, for 1~KiB, 64~KiB and 1~MiB bodies, since they
already stay on the worker. The inline executor with page-cache reads
switched off ran as fast: a read that the page cache answers does not
block, whichever call makes it.

\textbf{Unsynchronised stores} ran 24 to 32\,\% faster inline, with 4, 16
or 64 workers alike (Table~\ref{tab:executor-unsynced}): above the 20\,\%
that \reqref{HLR-SCOPE.V1-01} asks of a new mechanism, on the figure
\reqref{TC-PERFORMANCE.STORERATE-01} measures. With four workers, the
server's threads fell from about 90 to 8. More workers did not help the
pool: with 64, its rate stayed within 3\,\% of its rate with 4.

\begin{table}[tb]
\centering\small\renewcommand{\arraystretch}{1.3}%
\caption{Unsynchronised stores by executor: stores per second (median of
three rounds) and change against the pool with four workers; the 99th
percentile of latency (p99) and the server's threads for 4~KiB bodies}%
\label{tab:executor-unsynced}
\begin{tabular}{|l|r|r|r|r|}
\hline
\textbf{Executor, workers} & \textbf{4~KiB} & \textbf{64~KiB} &
\textbf{p99} & \textbf{Threads} \\
\hline
Pool, 4 & 8,883 & 7,158 & 12~ms & 88 \\
\hline
Pool, 64 & 9,092 ($+2$\,\%) & 6,990 ($-2$\,\%) & 36~ms & 385 \\
\hline
Inline, 4 & 11,612 ($+31$\,\%) & 8,948 ($+25$\,\%) & 10~ms & 8 \\
\hline
Inline, 16 & 11,652 ($+31$\,\%) & 8,991 ($+26$\,\%) & 21~ms & 20 \\
\hline
Inline, 64 & 11,723 ($+32$\,\%) & 8,889 ($+24$\,\%) & 42~ms & 68 \\
\hline
Block in place, 4 & 9,043 ($+2$\,\%) & 7,054 ($-1$\,\%) & 37~ms & 473 \\
\hline
\end{tabular}
\end{table}

\textbf{Synchronised stores} collapsed inline with few workers
(Table~\ref{tab:executor-strict}). A worker waiting in \texttt{fdatasync}
can do nothing else, so stores complete only as fast as the workers' syncs
do: 4 workers made 229 stores of 4~KiB a second, one per worker every
17.5~ms, and 16 workers made 802, one every 20~ms --- the time an
\texttt{fdatasync} takes on this host. With 64 workers, as many as there
were connections, the rate matched the pool's, which had grown to 150
threads and so never ran short of them, and the server used 10 to 14\,\%
less processor time than on the pool.

\begin{table}[tb]
\centering\small\renewcommand{\arraystretch}{1.3}%
\caption{Synchronised stores by executor: stores per second (median of
three rounds) and change against the pool with four workers; the 99th
percentile of latency (p99) and the processor time used, in processors, for
4~KiB bodies}%
\label{tab:executor-strict}
\begin{tabular}{|l|r|r|r|r|}
\hline
\textbf{Executor, workers} & \textbf{4~KiB} & \textbf{64~KiB} &
\textbf{p99} & \textbf{Processors} \\
\hline
Pool, 4 & 2,642 & 1,749 & 37~ms & 1.46 \\
\hline
Inline, 4 & 229 ($-91$\,\%) & 198 ($-89$\,\%) & 519~ms & 0.24 \\
\hline
Inline, 16 & 802 ($-70$\,\%) & 668 ($-62$\,\%) & 135~ms & 0.50 \\
\hline
Inline, 64 & 2,613 ($-1$\,\%) & 1,663 ($-5$\,\%) & 35~ms & 1.26 \\
\hline
Inline, 128 & 2,646 ($0$\,\%) & 1,720 ($-2$\,\%) & 34~ms & 1.28 \\
\hline
Block in place, 4 & 2,727 ($+3$\,\%) & 1,744 ($0$\,\%) & 33~ms & 1.53 \\
\hline
\end{tabular}
\end{table}

\textbf{More workers cost latency.} With 64 workers sharing four hardware
threads, the 99th percentile of 1~KiB hits rose from 2.3~ms to 14~ms, on the
pool as inline; with 16 workers, to 6.3~ms. A server that needs one worker
per concurrent sync to keep synchronised stores inline pays this on every
request.

\textbf{\texttt{block\_in\_place} saves nothing.} It ran as fast as the
pool, with a worse tail --- the 99th percentile of unsynchronised 4~KiB
stores rose from 12~ms to 37~ms --- and the server ran 473 threads. Before
each operation, it sends the worker's queue to a thread that it starts
through the blocking
pool:\footnote{\href{https://github.com/tokio-rs/tokio/blob/tokio-1.53.3/tokio/src/runtime/scheduler/multi_thread/worker.rs\#L477-L489}{Tokio
source (release 1.53.3):
\texttt{tokio/src/runtime/scheduler/multi\_thread/worker.rs}, lines
477--489}.} a handoff on every call, in another form.

\textbf{The test suite.} Inline, seven tests failed and two never ended.
The seven check that work leaves the worker, or the thread that drops a
handler: four of \reqref{TC-STORAGE.EXECUTOR-01}, and three others, for
aborted stores, purged records and published stores. The two that never
ended pause an update between its writes, to check that a lookup waits it
out (\reqref{TC-LAYOUT.OBJECTFILE-01}, \reqref{TC-STORAGE.RECLAIM-01}); the
test's runtime has one thread, which the paused update then holds. These
tests check \reqref{FDS-STORAGE.EXECUTOR-01}, which the inline executor
sets aside by design: a mode that runs work on the workers needs a
functional design and test cases of its own.

\section{Why a Write on the Worker Can Stall}
\label{app:executor-stalls}

The inline gain on unsynchronised stores rests on their writes not
blocking: a write to the page cache copies the bytes and returns. Two
mechanisms of the kernel can make it wait regardless, and on a worker that
wait stalls every task the worker holds.

\textbf{Dirty-page throttling.} Linux lets \emph{dirty} data --- written to
the page cache but not yet to the device --- grow up to a threshold, and
then makes the writers themselves wait for the device. Background
write-back starts at \texttt{vm.dirty\_background\_ratio} of the memory
available for the page cache, 10\,\% by default, and writers are throttled
once dirty data passes the midpoint between that and
\texttt{vm.dirty\_ratio}, 20\,\% by
default.\footnote{\href{https://www.kernel.org/doc/html/latest/admin-guide/sysctl/vm.html\#dirty-ratio}{Linux
kernel documentation: ``Documentation for /proc/sys/vm/''},
\texttt{dirty\_background\_ratio} and \texttt{dirty\_ratio}, both percentages
``of total available memory that contains free pages and reclaimable
pages''; the midpoint:
\href{https://github.com/torvalds/linux/blob/v6.12/mm/page-writeback.c\#L837-L841}{Linux
source (release 6.12): \texttt{mm/page-writeback.c}, lines 837--841}.} On this
host, at the time of writing, the kernel put the two thresholds at about
0.7 and 1.4~GiB, so writers would be throttled from about 1~GiB of dirty
data. A throttled writer sleeps for up to 200~ms at a time and, while dirty
data stays above the upper threshold, sleeps again; only a fatal signal
ends the sleep
early.\footnote{\href{https://github.com/torvalds/linux/blob/v6.12/mm/page-writeback.c\#L48}{Linux
source (release 6.12): \texttt{mm/page-writeback.c}, line 48}
(\texttt{MAX\_PAUSE}, a fifth of a second), and
\href{https://github.com/torvalds/linux/blob/v6.12/mm/page-writeback.c\#L2005-L2007}{lines
2005--2007} (the sleep, \texttt{TASK\_KILLABLE}).} The throttle applies to
every task that dirties pages while the threshold is exceeded. Under a burst
of writes --- fills of large objects arriving faster than the device writes
them --- every worker that writes stalls at once, and more workers help only
while some of them are not writing.

In a memory cgroup, as in most containers, the kernel applies the same
ratios to the memory the cgroup may use, and enforces whichever limit is the
more
restrictive:\footnote{\href{https://www.kernel.org/doc/html/latest/admin-guide/cgroup-v2.html\#writeback}{Linux
kernel documentation: ``Control Group v2''}, Writeback.} a container limited
to 2~GiB of memory throttles its writers from at most about 300~MiB of dirty
data.

Whether the store runs of Section~\ref{app:executor-experiment} reached the
throttle is not known. They may well have: the 64~KiB runs wrote bodies at
about 560~MiB a second for 10~seconds, several times the threshold, but the
script recorded no dirty-data figures.

\textbf{The ext4 journal.} On ext4, a write that changes a file's timestamps
marks its inode dirty, which starts a journal handle, and so does every
create and
rename.\footnote{\href{https://github.com/torvalds/linux/blob/v6.12/fs/ext4/inode.c\#L5977-L5981}{Linux
source (release 6.12): \texttt{fs/ext4/inode.c}, lines 5977--5981}
(\texttt{ext4\_dirty\_inode}).} While the journal commits its running
transaction, starting a handle waits for the commit to end, and that wait
cannot be
interrupted.\footnote{\href{https://github.com/torvalds/linux/blob/v6.12/fs/jbd2/transaction.c\#L142-L150}{Linux
source (release 6.12): \texttt{fs/jbd2/transaction.c}, lines 142--150}
(\texttt{wait\_transaction\_locked}, \texttt{TASK\_UNINTERRUPTIBLE}).}
Commits follow \texttt{fdatasync} calls, such as synchronised stores make,
and the journal's periodic commit, every 5~seconds by
default;\footnote{\href{https://www.kernel.org/doc/html/latest/admin-guide/ext4.html}{Linux
kernel documentation: ``ext4 General Information''},
\texttt{commit=nrsec}.} in the default \texttt{data=ordered} mode, a commit
first writes out the data of the files it covers. A synchronised store on
one worker can thus hold up an unsynchronised store's write on another for
as long as the device takes.

\textbf{No non-blocking writes on ext4.} \texttt{preadv2} with
\texttt{RWF\_NOWAIT} lets a read fail rather than wait
(\reqref{ADR-STORAGE.PAGECACHE-01}); \texttt{pwritev2} takes the same flag,
although its manual page documents it for reads
only.\footnote{\href{https://man7.org/linux/man-pages/man2/pwritev2.2.html}{Linux
manual page: \texttt{pwritev2(2)}}: ``Currently, this flag is meaningful
only for preadv2().''} In Linux 6.12, ext4 refuses it for buffered writes,
with
\texttt{EOPNOTSUPP},\footnote{\href{https://github.com/torvalds/linux/blob/v6.12/fs/ext4/file.c\#L291-L292}{Linux
source (release 6.12): \texttt{fs/ext4/file.c}, lines 291--292}.} as this host
confirmed; so does tmpfs. On ext4, then, a write on the worker can neither
be kept from stalling nor be cut short once it has.
Section~\ref{app:executor-fs} lists the file systems that accept the flag.

\section{Directions for a Later Version}
\label{app:executor-directions}

\subsection{Bounded inline writes}
\label{app:executor-bounded}

Of the directions in this section, this is the one the investigation
favours, as a mode chosen by configuration beside today's. The backend
keeps a fixed number $K$ of
\emph{inline slots}, fewer than the workers. The operations of an
unsynchronised store that do not wait for the device by nature --- creating
its temporary, writing it, renaming it --- take a slot, with one atomic
operation, and run on the worker; when every slot is taken, they go to the
pool, as today. \texttt{fdatasync}, and the reads that the page cache cannot
answer, always go to the pool.
\begin{itemize}
  \item In normal operation, writes take microseconds and the slots are
    rarely all taken: most stores make no handoff, and keep most of the gain
    measured above.
  \item Under a burst of writes, at most $K$ workers stall. The others go on
    serving hits, and further writes queue on the pool. The server slows
    down during the burst, but does not stop.
  \item The cost falls on the tasks that a stalled worker holds: they wait
    out the stall, up to 200~ms for each throttle pause. That includes the
    task in the worker's \emph{LIFO slot}, the one it is to poll next, which
    other workers cannot
    steal.\footnote{\href{https://github.com/tokio-rs/tokio/blob/tokio-1.53.3/tokio/src/runtime/builder.rs\#L1310-L1313}{Tokio
    source (release 1.53.3): \texttt{tokio/src/runtime/builder.rs}, lines
    1310--1313}: ``this slot cannot be stolen by other worker threads''.}
  \item With $K = 0$, the backend behaves as this version does.
\end{itemize}
Open questions:
\begin{itemize}
  \item The backend does not see \texttt{ServerConf::threads}, so it cannot
    check that $K$ is below the number of workers. The option could take
    both numbers, or the backend could read the number of workers from the
    runtime, once.
  \item With work stealing off, each Pingora worker runs a runtime of its
    own, and nothing takes over the tasks of a stalled worker: the mode
    needs work stealing on.
  \item The deviation \reqref{ADR-STORAGE.EXECUTOR-01} records --- an
    upstream \texttt{103 Early Hints} lost while a sole worker was blocked
    --- came from a stalled worker: \reqref{TC-COMPLIANCE.RFC-01}'s
    differential would have to pass in this mode too.
  \item \reqref{FDS-STORAGE.EXECUTOR-01} and its test cases describe the mode
    with $K = 0$; the others need a functional design and test cases of their
    own.
\end{itemize}

\subsection{Early write-back}
\label{app:executor-writeback}

After publishing an unsynchronised store, the backend could ask the kernel
to start writing its data to the device, without waiting for it:
\texttt{sync\_file\_range} with \texttt{SYNC\_FILE\_RANGE\_WRITE}. The call
can still block when the device's queue is
full,\footnote{\href{https://man7.org/linux/man-pages/man2/sync_file_range.2.html}{Linux
manual page: \texttt{sync\_file\_range(2)}}: ``Note that even this may block
if you attempt to write more than request queue size.''} so it would run on
the pool, with nothing waiting for it. The cache's own dirty data would then
stay small, so the cache alone would rarely push the host towards the
throttle, and a host crash would lose fewer stores. The cost: the device
writes each store sooner, including stores that are evicted or replaced
before the kernel would have written them, and a store makes one more
handoff. It does nothing against the dirty data that other programs on the
host make.

\subsection{Non-blocking writes where the file system has them}
\label{app:executor-nowait-writes}

Writes could follow \reqref{ADR-STORAGE.PAGECACHE-01}: write on the worker
with \texttt{pwritev2} and \texttt{RWF\_NOWAIT}, and run the store on the
pool from the first call that reports it would wait. On XFS, such a write
returns \texttt{EAGAIN} instead of waiting for the file's I/O
lock,\footnote{\href{https://github.com/torvalds/linux/blob/v6.12/fs/xfs/xfs_file.c\#L194-L196}{Linux
source (release 6.12): \texttt{fs/xfs/xfs\_file.c}, lines 194--196}.} for the
dirty-page
throttle,\footnote{\href{https://github.com/torvalds/linux/blob/v6.12/fs/iomap/buffered-io.c\#L918}{Linux
source (release 6.12): \texttt{fs/iomap/buffered-io.c}, line 918} and
\href{https://github.com/torvalds/linux/blob/v6.12/fs/iomap/buffered-io.c\#L932-L933}{lines
932--933}; \href{https://github.com/torvalds/linux/blob/v6.12/mm/page-writeback.c\#L2001-L2003}{\texttt{mm/page-writeback.c},
lines 2001--2003}.} or to update the file's
timestamps.\footnote{\href{https://github.com/torvalds/linux/blob/v6.12/fs/inode.c\#L2310-L2314}{Linux
source (release 6.12): \texttt{fs/inode.c}, lines 2310--2314}.} The last will
matter most: a write that would change the file's modification time fails,
and on a file being written that is any write made after the clock has
moved since the previous one. How often stores would fall back to the pool
needs measuring. Creating the temporary and renaming it stay blocking calls
on every file system: \texttt{openat2} refuses \texttt{RESOLVE\_CACHED} for
a
create,\footnote{\href{https://github.com/torvalds/linux/blob/v6.12/fs/open.c\#L1331-L1334}{Linux
source (release 6.12): \texttt{fs/open.c}, lines 1331--1334}.} and no flag makes a
rename fail rather than wait.

As with page-cache reads, this would be an option, checked once when the
cache is constructed --- a file system that refuses the flag then makes the
construction fail, naming the option --- and never a silent detour at run
time.

\subsection{io\_uring}
\label{app:executor-iouring}

\texttt{io\_uring} submits operations through a ring of memory shared with
the kernel, and reports their completions through another, with no thread
of the process waiting on them. The kernel first tries each operation
without blocking, and hands one that would block to worker threads of its
own. What that means for the backend's operations, in Linux 6.12:
\begin{itemize}
  \item A read that misses the page cache needs no thread on ext4, XFS or
    btrfs: the kernel completes it when the device delivers the
    data.\footnote{\href{https://github.com/torvalds/linux/blob/v6.12/io_uring/rw.c\#L752-L753}{Linux
    source (release 6.12): \texttt{io\_uring/rw.c}, lines 752--753}: the file
    system must declare \texttt{FOP\_BUFFER\_RASYNC}, as those three do.}
    This is the case \reqref{HLR-PERSISTENCE.CAPACITY-01} makes common, and
    today it costs a failed page-cache read and a handoff.
  \item A buffered write runs without blocking on XFS, as above. On ext4, it
    always goes to the kernel's
    threads.\footnote{\href{https://github.com/torvalds/linux/blob/v6.12/io_uring/rw.c\#L1055-L1059}{Linux
    source (release 6.12): \texttt{io\_uring/rw.c}, lines 1055--1059}.}
  \item \texttt{fdatasync}, \texttt{sync\_file\_range}, renames and unlinks
    always go to the kernel's
    threads.\footnote{\href{https://github.com/torvalds/linux/blob/v6.12/io_uring/sync.c\#L76-L77}{Linux
    source (release 6.12): \texttt{io\_uring/sync.c}, lines 76--77} and
    \href{https://github.com/torvalds/linux/blob/v6.12/io_uring/sync.c\#L45-L46}{45--46};
    \href{https://github.com/torvalds/linux/blob/v6.12/io_uring/fs.c\#L77}{\texttt{io\_uring/fs.c},
    line 77} and \href{https://github.com/torvalds/linux/blob/v6.12/io_uring/fs.c\#L126}{126}.}
\end{itemize}
On ext4, then, \texttt{io\_uring} would move stores from Tokio's pool to the
kernel's, at a cost not measured; its clear gain is on reads from the
device.

Integration is the larger task, and Tokio offers two starting points.
\texttt{tokio-uring} drives \texttt{io\_uring} from a runtime of its
own,\footnote{\href{https://github.com/tokio-rs/tokio-uring\#getting-started}{\texttt{tokio-uring}
README}: ``Using \texttt{tokio-uring} requires starting a
\texttt{tokio-uring} runtime.''} which Pingora does not use. Tokio itself
has experimental support: built with \texttt{-{}-cfg tokio\_unstable} and
its \texttt{io-uring} feature, a runtime that enables all its drivers ---
as the Pingora runtimes
do\footnote{\texttt{pingora-runtime} calls
\pgsrc{pingora-runtime/src/lib.rs}{470}{\texttt{enable\_all}} on a
work-stealing runtime (\pgsrc{pingora-runtime/src/lib.rs}{627}{L627} on the
others).} --- also drives a ring, which \texttt{tokio::fs} uses
to open, read, write and rename
files.\footnote{\href{https://github.com/tokio-rs/tokio/blob/tokio-1.53.3/tokio/src/lib.rs\#L481-L482}{Tokio
source (release 1.53.3): \texttt{tokio/src/lib.rs}, lines 481--482}, and
\href{https://github.com/tokio-rs/tokio/blob/tokio-1.53.3/tokio/src/runtime/builder.rs\#L384-L391}{\texttt{tokio/src/runtime/builder.rs},
lines 384--391}.} It offers no positional reads or writes, which this
layout depends on, and the opt-in to unstable features belongs to the
integrator's build, not to this backend. Short of either, the backend would
drive a ring itself, through the \texttt{io-uring} crate, and wake each task
when its completion arrives. Two constraints follow:
\begin{itemize}
  \item The kernel owns a buffer until the operation that uses it
    completes. A handler that Pingora drops --- its client gone --- cannot
    free its buffer before then: buffers must belong to the operation in
    flight, and an abandoned operation must be cancelled or awaited.
  \item Every operation the backend makes needs its own opcode: renames and
    unlinks exist since Linux 5.11, \texttt{openat2} and plain writes since
    5.6,\footnote{\href{https://man7.org/linux/man-pages/man2/io_uring_enter.2.html}{Linux
    manual page: \texttt{io\_uring\_enter(2)}}, \texttt{IORING\_OP\_RENAMEAT},
    \texttt{IORING\_OP\_UNLINKAT}, \texttt{IORING\_OP\_OPENAT2},
    \texttt{IORING\_OP\_WRITE}.} so the mode would need Linux 5.11 or later.
\end{itemize}

Deployment limits it further. The default seccomp profile of Docker Engine
has blocked the \texttt{io\_uring} system calls since release~25, so in a
default container \texttt{io\_uring\_setup}
fails.\footnote{\href{https://github.com/moby/moby/pull/46762}{moby/moby
pull request 46762}, ``seccomp: block io\_uring\_* syscalls in default
profile'', merged on 2~November~2023: \texttt{profiles/seccomp/default.json}
allows the calls at tag \texttt{v24.0.0} and no longer at \texttt{v25.0.0}.}
Google disabled it on its production servers after 60\,\% of the
submissions to its kernel exploit reward programme, in the year to
June~2023, used
it,\footnote{\href{https://security.googleblog.com/2023/06/learnings-from-kctf-vrps-42-linux.html}{Google
Security Blog, ``Learnings from kCTF VRP's 42 Linux kernel exploits
submissions''}, 14~June~2023.} and Linux 6.6 added a setting that disables
it for the whole
host.\footnote{\href{https://www.kernel.org/doc/html/latest/admin-guide/sysctl/kernel.html\#io-uring-disabled}{Linux
kernel documentation: ``Documentation for /proc/sys/kernel/''},
\texttt{io\_uring\_disabled};
\href{https://kernelnewbies.org/Linux_6.6}{Linux 6.6 release notes
(kernelnewbies.org)}.} \texttt{io\_uring} can therefore only be an option
beside the pool, checked when the cache is constructed, as page-cache reads
are. It is the largest of the mechanisms in this appendix, and
\reqref{HLR-SCOPE.V1-01} would weigh it accordingly.

\subsection{Not pursued}
\label{app:executor-not-pursued}

\begin{itemize}
  \item \emph{A lazy handoff.} The Go runtime lets a goroutine make a
    blocking system call on its own thread. A monitor thread, which checks
    every 20~microseconds while it finds work to do and backs off to every
    10~milliseconds while it finds none, hands the thread's queue of
    goroutines to another thread if the call is still running at its next
    check.\footnote{\href{https://github.com/golang/go/blob/go1.25.0/src/runtime/proc.go\#L6226-L6233}{Go
    source (release 1.25.0): \texttt{src/runtime/proc.go}, lines
    6226--6233} (the monitor's period),
    \href{https://github.com/golang/go/blob/go1.25.0/src/runtime/proc.go\#L6413}{line
    6413} and
    \href{https://github.com/golang/go/blob/go1.25.0/src/runtime/proc.go\#L6441}{line
    6441} (the handoff).} A short call costs nothing, and a long one a
    single handoff. Tokio has no such monitor, and adding one is a change to
    Tokio, out of this backend's reach.
  \item \emph{Routing writes by the host's dirty data}, read from
    \texttt{/proc/vmstat} before each store. It would detect the host's
    state on the hot path, which this project avoids: its features are
    chosen by configuration and checked once, as page-cache reads are. The
    reading would also be stale by the time the write runs.
\end{itemize}

\section{File Systems for the Cache}
\label{app:executor-fs}

What the investigation suggests for the file system that holds the cache
directory. Only ext4 has been measured with this backend; a later version
should confirm the rest.
\begin{itemize}
  \item \emph{XFS}, on Linux 6.0 or later, is the one to evaluate first. Of
    ext4, XFS and btrfs, it is the only one whose ordinary files accept
    non-blocking buffered
    writes,\footnote{\href{https://github.com/torvalds/linux/blob/v6.12/fs/xfs/xfs_file.c\#L1597-L1598}{Linux
    source (release 6.12): \texttt{fs/xfs/xfs\_file.c}, lines 1597--1598}
    (\texttt{FOP\_BUFFER\_WASYNC}); since Linux 6.0, per its
    \href{https://kernelnewbies.org/Linux_6.0}{release notes
    (kernelnewbies.org)}.} which the direction of
    Section~\ref{app:executor-nowait-writes} needs, and page-cache reads
    should work on it as on ext4. How it compares with ext4 on synchronised
    stores, which their \texttt{fdatasync} dominates, is not measured.
  \item \emph{ext4}, on which every figure in this specification was
    measured, serves both kinds of store, with page-cache reads. A write on
    it cannot avoid the waits of Section~\ref{app:executor-stalls}. When
    stores are unsynchronised, keep its default \texttt{data=ordered} mode
    (\reqref{ADR-PERSISTENCE.UNSYNCED-01}).
  \item \emph{btrfs} accepts non-blocking buffered writes only for files
    that do not use copy-on-write; for the others, the default, it returns
    \texttt{EAGAIN} at
    once.\footnote{\href{https://github.com/torvalds/linux/blob/v6.12/fs/btrfs/file.c\#L1159-L1161}{Linux
    source (release 6.12): \texttt{fs/btrfs/file.c}, lines 1159--1161}.} A file
    without copy-on-write also has no data
    checksums,\footnote{\href{https://btrfs.readthedocs.io/en/latest/Administration.html}{btrfs
    documentation: ``Administration''}, mount options: ``Nodatacow implies
    nodatasum, and disables compression.''} which are the reason to choose
    btrfs (\reqref{ADR-PERSISTENCE.DURABILITY-01}). It serves today's modes,
    and offers nothing to the directions above.
  \item \emph{tmpfs} does not suit: it refuses \texttt{RWF\_NOWAIT} for
    reads and writes alike, so page-cache reads cannot be used on it; it
    loses its contents when the host restarts; and it holds no more than
    memory, against \reqref{HLR-PERSISTENCE.CAPACITY-01}.
  \item \emph{Network file systems} are ruled out already: the cache
    directory must be on a local file system
    (\reqref{ADR-STORAGE.WRITERID-01}).
\end{itemize}

Two choices help whatever the file system:
\begin{itemize}
  \item \emph{A file system of its own.} On a file system it does not
    share, the cache's stores wait on no other program's journal commits,
    and other programs wait on none of the cache's. The dirty-data limits
    remain the host's, or the container's, so throttling stays shared.
  \item \emph{\texttt{noatime}.} Under \texttt{relatime}, the Linux default,
    the first read of a record after it is written updates its access time.
    On ext4 that is a journalled inode update, and a page-cache read
    (\reqref{ADR-STORAGE.PAGECACHE-01}) makes it on the worker. Mounting the
    file system \texttt{noatime} removes those updates.
\end{itemize}

\section{What a Later Version Should Measure}
\label{app:executor-measure}

\begin{itemize}
  \item \emph{Hits during a burst of writes}: the latency of hits while
    unsynchronised stores of large objects keep dirty data above the
    throttle, with the pool, with bounded inline writes for several values
    of $K$, and with non-blocking writes on XFS; on the host, and in a
    container with a memory limit; recording the dirty and write-back
    figures of \texttt{/proc/meminfo} throughout. The script used so far
    runs one load at a time, and would need a second load generator.
  \item \emph{Reads from the device}: every hit measured so far was on an
    object the page cache held. Reads that miss it, the normal case for a
    cache larger than memory, should be measured on the pool and, if it is
    prototyped, with \texttt{io\_uring}.
  \item \emph{Synchronised stores on XFS}, against ext4 on the same device.
  \item \emph{Early write-back}: store rate, dirty data and device writes,
    with and without it.
\end{itemize}

No performance requirement of this version bounds the latency of hits while
stores run: \reqref{HLR-PERFORMANCE.LOOKUP-01} bounds lookup latency against
the size of the cache, and \reqref{HLR-PERFORMANCE.HITRATE-01} and
\reqref{HLR-PERFORMANCE.STORERATE-01} bound rates. A later version that keeps
writes on the workers would need one, so that the gain on stores, which
\reqref{HLR-SCOPE.V1-01} counts, is weighed against the stalls it risks.
